\section{Community and Society}

As I was working on a post, three items came to my attention for review. Since they touch on similar topics, I'll weave traditional notions as part of the review. These are:

\begin{enumerate}


\item \emph{From the German Conservative Revolution to the New Right} (200 pages, softcover), by \textbf{Lucian Tudor}. The book deals with writers such as Arthur Moeller van den Bruck, Oswald Spengler, Othmar Spann, etc.

It is published by Identitas/Círculo de Investigaciones PanCriollistas\footnote{\url{http://fni.cl/identitas}}, an identitarian group of Latin American Creoles.

\item \emph{Transcendence \& the Aristocratic Principle}, by \textbf{Edwin Dyga}, published in Aristokratia III\footnote{\url{http://www.aristokratia.info/buy-aristokratia.html}}

\item \emph{Testament de Ybarzabal}, an obscure work by \textbf{Sidney Vigneaux}, which \textbf{Rene Guenon} considered to be of interest from this review of the Protocols of Zion\footnote{\url{http://www.hyperion-journal.net/on-the-protocols-of-zion.html}}. (H/T Avery\footnote{\url{http://avery.morrow.name/blog/}} for pointing this out and locating the French text.) This will come later, since I have to translate the French text while Aeneas is translating the Evola essay mentioned by Guenon.
\end{enumerate}


The first two works will be reviewed from a specific perspective, which we will make clear. Unfortunately, there is the persistent idea that the New Right and Tradition have something in common and are aiming for the same goal. It is clear from these works that the New Right is a patchwork of ideas with little in common among those associated, either willingly or unwillingly, with that movement. Fundamental terms such as ``Right'', ``Identity'', and the ``modern world'' are never clearly defined, certainly not in a positive way.

Unlike the New Right, Tradition is not a ``movement'' among other movements vying for power or influence. Rather, Tradition sees itself as the True, which supersedes all ``movements''.

So let's start with some premises as a working hypothesis. Of course, since Tradition is understood in its fullness only by a higher intellect, there is not, and can never be, scientific proof or a logical demonstration of its truthfulness. Nevertheless, it is not irrational so these premises can certainly be grasped with a little effort. If you disagree, or are uncertain, treat it as a ``what if'', or a language game that you can learn to play once you figure out the rules. If you think about it, you may even come to see that they are true.

\paragraph{Tradition}


Great civilizations have been founded on certain metaphysical principles which Rene Guenon has brought to light. These principles he called ``Tradition''. However, that is an abstract term, a genus, that requires a species. Hence, you can no more be a ``Traditionalist'' than you can paint your house a ``color''. There is always the follow up question, ``Which color?'', ``Which Tradition?'' Qualifying Tradition with the adjective ``radical'' adds nothing although it sounds way cool in some circles.

What, then, is the Tradition of the West? Here we agree with St Augustine's assertion that there has been but one Tradition, now known as Catholic. In this, Guenon agrees, and its fulfillment was the civilization of the European Middle Ages. Hence, to follow Tradition in the West means to be a Catholic (with some qualifications to follow). If Tradition is true, then a fortiori the Western tradition is true.

\paragraph{Degeneration and Regeneration of Tradition}


Now Guenon was motivated to a large extent by what he saw as the decline of Tradition in the West. A Tradition in itself cannot degenerate. Rather it can be forgotten, neglected, or the titular heads of its institution fail to understand and transmit it.

This is a point in which we agree with the Association Boris Mouravieff\footnote{\url{http://www.association-boris-mouravieff.com/pages/ConTra.html}}:

\begin{quotex}
If Tradition, and the teachings attached to it, follow the historic evolution of Man and Nature since Creation, it doesn't generally appear in plain sight. Esoteric teaching has always been given to certain persons most often via oral transmission. Christian esoteric tradition, notably, has been preserved intact for centuries under the protection of Hermetism. ... Hermetism has constituted for quite some time the safeguard for tradition.
\end{quotex}


As \textbf{Valentin Tomberg} describes it, Hermetism has followed a parallel path to the Institution, preserving its teachings independently while not actively opposing it. The time has come to make these teachings more public, although oral transmission can still not be dispensed with. Hence, group work is actually the most important part of Gornahoor.

\paragraph{Community and the Individual}


The importance of community and the rejection of individualism is common to both Tradition and the New Right. Unfortunately, we have poor examples to follow. Guenon became a Sufi as part of his ``personal equation''. Evola created a post-Christian idiosyncratic personal philosophy.

Guenon justified his choice on the grounds that some men are advanced to be able to choose their own tradition. In a sense, that is true. Tomberg claims in Letter XX that such an individual is less tied to the family, people or race. (575) However, the higher individual is born into a community that is suitable for him:

\begin{quotex}
The individuality, in the case where his incarnation is ruled by the law of the vertical, descends consciously of his own free will to birth, into an environment where he is wanted and awaited.
\end{quotex}


So it would be unusual to reject that environment, which includes family, nation, race, religious tradition and so on. In other words, while most men may be attached to a community by family, habit, custom, or social pressure, the conscious individual chooses to belong to it by his own free will. In other words, his soul elements are in harmony, \emph{viz}., his instinctual attachments, emotional bonds, intellectual commitment, and will all align. That is \textbf{Identity}.

\paragraph{Loyalty and Fidelity}


\begin{quotex}
Faithless is he that says farewell when the road darkens.
\flright{\textsc{J.R.R.Tolkien}}
\end{quotex}


Tudor, in his book, mentions Ferdinand Tonnies' distinction between Community and Society. Community consists of the organic relations and a sense of connection and belonging which arise as a result of natural will, while Society consists of mechanical or instrumental relations which are consciously established and thus the result of rational will.

Community designates a social entity which is based upon solidarity, bonding, a sense of connectedness and interdependence; it means belonging to a supra-individual whole on a deep spiritual level. In a Society, the unity between individuals exists only on a superficial level, where individuals are not united by spiritual bonds but by practical, legal, economic, or other artificial relationships.

Keep this useful distinction in mind. At the end we will see that Tradition is a Community and the New Right a Society. It is Gornahoor's contention that the organization of an intellectual movement should mimic its doctrine.

{Next: paganism, the right, and identity}


\flrightit{Posted on 2015-05-21 by Cologero}

\begin{center}* * *\end{center}

\begin{footnotesize}
\begin{sffamily}
\texttt{White Rabbit on 2015-05-22 at 20:11 said:}

`` I t would be unusual to reject that environment, which includes family, nation, race, religious tradition and so on.''

When those elements don't coincide, it seems like a conundrum.

\hfill

\texttt{Lucian Tudor on 2015-05-26 at 20:43 said:}

I would like to make a brief remark about this statement near the end of the post: ``At the end we will see that Tradition is a Community and the New Right a Society.'' What I can certainly agree with is that all true traditional societies are based upon the value of community and supra-individualist (thus anti-individualist) principles. However, the quoted statement is rather confusing, as I am not sure how the concepts of gemeinschaft (community) and gesellschaft (society) can be applied to the New Right (a broad ideological group consisting of a variety of movements in a variety of nations) or Tradition (a set of perennial metaphysical principles in certain philosophical lines). Also, while the New Right as an international movement is -- using the term very loosely -- something of a large-scale gesellschaft, on the smaller scales it seeks and successfully builds true communities. For example, GRECE constructs communities at the various towns and localities where it sets its circles/centers; members and associates affiliate with each other to actively build communities with true solidarity.

\hfill

\texttt{Lucian Tudor on 2015-05-26 at 20:45 said:}

Thank you to Cologero, by the way, for discussing and referencing my book. Overall, the discussion on community, identity, etc. is rather agreeable to me and matches the positions I argued for in my own writings.

\hfill

\texttt{Cologero on 2015-05-27 at 07:48 said:}

Mr Tudor, you've made my point about ``Tradition''. As we've emphasized here, over and over, you must follow a particular Tradition, there is not ``tradition'' as such since it is an abstraction. So yes, a particular tradition is a community in that sense.

I am more interested in these GRECE-constructed communities. Actual examples would be of interest and would show the New Right to be more than an effete intellectual game. For example, what is the actual basis of the solidarity? Obviously, a ``wide variety of movements'' can only be a moving target, not the idee-force of an identity.

How are such communities organized? Obviously, they must be anti-democratic and anti-liberal. How, then, is the leadership chosen? Do they have a positive birthrate? Who decides who the members are?

\hfill

\texttt{Lucian Tudor on 2015-05-27 at 09:13 said:}

Yes, the New Right in Europe is far from being composed of pure intellectualism, even if its foundational activity is ``Right-wing'' intellectual activism. The GRECE communities are founded on the idea of establishing a deep sense of gemeinschaft (as opposed to just gesellschaft), as one can see it in close-knit villages, religious organizations, or large groups of friends. The various members who are close enough meet with each other do so not merely at gatherings devoted to intellectual activity (conferences, lectures, discussions, etc.), but also for purely social activities, to celebrate holidays, important personal events of a member (marriages, etc.), and even religious rituals in some cases. Not all members are friends with each other, but all engage in these activities as proper traditional communities do, which establishes the organic solidarity. I'm afraid I don't have all the details to answer your questions fully, but my suggestion would be to ask GRECE members or to read any works they have written on these matters (probably mostly in French). I am aware that this has been discussed to some extent Alain de Benoist's book ``Mémoire vive'' (or ``Mein Leben'' in German), but there is surely more information also in some other writings of his.

\end{sffamily}
\end{footnotesize}

